% =====================================================================
%  Chapter 21 opener -- "Chapter at a Glance", three reading paths + the
%  ladder staircase, THE running examples (the 16-slot array A shared with
%  Ch.18/Ch.19/Ch.20, now read through a MERGE-SORT tree -- every node keeps
%  its interval's SORTED list -- and through ONE persistence act A[3]=9 that
%  path-copies exactly 5 of the 31 sum-tree nodes, cross-checkable against
%  Ch.19's grid; plus the 10-node tree T of Ch.12-17 with its Ch.16 preorder
%  tin/tout and a weight assignment for the 21.4 four-root path trace),
%  the merge-sort / persistence / chairman / sparse taxonomy, the
%  who-owns-what comparison (Ch.19 owns the plain tree; Ch.20 owns the lazy
%  tag; Ch.11 owns the persistent TRIE; Ch.5 owns flattening; Ch.22 owns HLD;
%  Ch.23 owns beats; THIS chapter owns merge-sort trees, whole-structure
%  persistence, prefix-version k-th, and the sparse-table RMQ), recognition
%  guide, constraint->tool table, notation.
%  Exports fig:c19-mst and fig:c19-taxonomy.
%  Hand-computed values the whole chapter re-uses (verified once here).
%  Running array (16 slots, shared with Ch.18/Ch.19/Ch.20):
%    A = (3,2,4,5,1,1,5,3,7,0,2,4,1,6,3,2),  sum A = 49.
%  Sum segment tree, recursive heap layout (root t[1], children t[2v],
%  t[2v+1]; n=16 is a power of two so the tree is complete, 31 nodes, and
%  the leaves t[16..31] = A[1..16]).  The base internal grid t[1..15]:
%    node:   t1   t2   t3   t4  t5  t6  t7  t8 t9 t10 t11 t12 t13 t14 t15
%    intvl: 1-16 1-8  9-16 1-4 5-8 9-12 13-16 1-2 3-4 5-6 7-8 9-10 11-12 13-14 15-16
%    SUM:   49   24   25   14  10  13  12   5  9   2   8   7   6   7   5
%  MERGE-SORT overlay: the SAME heap layout, each node now stores its
%  interval's SORTED list (only the top three levels are listed here; every
%  deeper node is a sub-list of the one above it):
%    t1 [1-16] = (0,1,1,1,2,2,2,3,3,3,4,4,5,5,6,7)   (the whole array, sorted)
%    t2 [1-8]  = (1,1,2,3,3,4,5,5)      t3 [9-16] = (0,1,2,2,3,4,6,7)
%    t4 [1-4]  = (2,3,4,5)  t5 [5-8] = (1,1,3,5)   t6 [9-12] = (0,2,4,7)  t7 [13-16] = (1,2,3,6)
%  count<=3 on [2..9] (canonical nodes t17=(2),t9=(4,5),t5=(1,1,3,5),leaf9=(7)):
%    1 + 0 + 3 + 0 = 4.   (A[2..9]=(2,4,5,1,1,5,3,7); the <=3 values are 2,1,1,3.)
%  PERSISTENCE overlay (the one hand-applied write the whole chapter reuses):
%  version v0 is the base sum tree; the write A[3]=9 (was 4, delta +5) copies
%  ONLY the root-to-leaf-18 path -- t1,t2,t4,t9,t18 (5 of 31 nodes) -- giving
%  new sums t1=54, t2=29, t4=19, t9=14, t18=9; every other node is SHARED with
%  v0 untouched.
%  Shared tree T of Ch.12-17 (edges 1-2,1-3,2-4,2-5,3-6,3-7,5-8,7-9,9-10),
%  Ch.16 preorder tin=(1,2,6,3,4,7,8,5,9,10), INCLUSIVE tout=(10,5,10,3,5,
%  7,10,5,10,10): subtree(2)=[2,5]={2,4,5,8}, subtree(7)=[8,10]={7,9,10}.
%  21.4 weights w(1..10) = (3,2,4,5,1,1,5,3,7,0): path 3-8 = {3,1,2,5,8}
%  weights (4,3,2,1,3) sorted (1,2,3,3,4); path 8-10 = {8,5,2,1,3,7,9,10}
%  weights (3,1,2,3,4,5,7,0) sorted (0,1,2,3,3,4,5,7), so kth(8-10,k=6)=4.
%  Every overlay below is hand-computed; every later section re-uses exactly
%  these numbers -- if a section disagrees with this opener, the OPENER is
%  authoritative: fix the section, never the other way.
% =====================================================================

\begin{chaptertheorybox}{Chapter at a Glance}
Chapters~18, 19 and~20 all keep \emph{one} live structure and answer questions
about its \emph{present} state: the Fenwick tree folds sums, Ch.19's segment
tree folds any associative aggregate over an interval, Ch.20 bolts on a lazy
tag to update whole ranges.  This chapter is the four questions none of them
can answer, and the structures that can.  A \term{merge-sort tree} puts each
node's interval \emph{sorted list} beside Ch.19's single aggregate, so a range
can be counted \emph{by value} -- "how many of \(a[l..r]\) are \(\le x\)",
"what is the \(k\)-th smallest of \(a[l..r]\)" -- on an array that never
changes.  \term{Persistence by path copying} refuses to overwrite: a write
copies only the \(O(\log n)\) nodes on the root-to-leaf path and \emph{shares}
the rest, so every old version survives and you can read the past.  Folding
persistence into the value axis gives the \term{chairman} (prefix-version) tree
-- version \(i\) is the frequency histogram of \(a[1..i]\), and a \(k\)-th
smallest in \([l,r]\) is a joint descent on two roots in \(\bigO{\log V}\).
Finally the \term{sparse table} answers static idempotent range min/max/gcd in
\(O(1)\) by \emph{overlapping} two pre-computed blocks -- the faster structure
Chapters~19 and~20 kept promising but never built.

\textbf{You will learn to}
\begin{itemize}[nosep]
  \item build a merge-sort tree (sorted list per node, \(\bigO{n\log n}\)) and
        count values \(\le x\) in a range in \(\bigO{\log^2 n}\), or \(k\)-th
        by an external binary search (\(\secref{c19-mst}\), SPOJ~KQUERY); \hfill\lv{2}--\lv{3}
  \item \emph{path-copy}: keep a root per version, copy only the touched path
        into a node pool, read any past version, and budget
        \(n+q\log n\) nodes so the pool never overflows
        (\(\secref{c19-copy}\), CF~707D); \hfill\lv{3}--\lv{4}
  \item build prefix-version (chairman) trees and answer the \(k\)-th smallest
        in \([l,r]\) by descending two roots at once (\(\secref{c19-chairman}\),
        SPOJ~MKTHNUM); \hfill\lv{4}
  \item plant a persistent root at every DFS node of a tree and answer a
        \(k\)-th / count on a \(u\)--\(v\) \emph{path} from FOUR roots plus
        LCA (\(\secref{c19-ontrees}\), SPOJ~COT); \hfill\lv{4}--\lv{5}
  \item build a sparse table for \(O(1)\) static idempotent RMQ, and \emph{choose}
        among all these engines on sight of the constraints
        (\(\secref{c19-sparse}\)). \hfill\lv{3}--\lv{5}
\end{itemize}

\textbf{Prerequisites.} Chapter~19 (\chref{c13}) is assumed -- the heap-layout
segment tree whose interval geometry every figure here reuses.  Chapter~20
(\chref{c15}) for the \emph{contrast} in \(\secref{c19-copy}\): a lazy tag is
precisely the thing persistence refuses to share.  Chapter~11's tries
(\chref{c08}) already built a persistent \emph{trie} (\defref{c08-persistent});
this chapter generalises that idea to whole structures and never re-builds the
trie.  Chapter~11 LCA (\chref{c11}) for \(\mathit{lca}()/\)binary lifting and
Chapter~5 (\chref{c24}) for \(\tin/\tout\), both cited by label in
\(\secref{c19-ontrees}\).
\end{chaptertheorybox}

% ---------------- Reading paths + ladder position ----------------
\begin{chaptertheorybox}{How to Read This Chapter (three paths)}
Level badges \(\lv2\)--\(\lv5\): \(\lv2\) is reading a sorted list off a drawn
merge-sort node and bisecting it; \(\lv3\) is a working merge-sort-tree count
on a real input (SPOJ~KQUERY is exactly this) or a hand-built path-copy of one
leaf; \(\lv4\) is a node-pool persistence with many versions, or the chairman
two-root \(k\)-th descent, or planting roots on a tree; \(\lv5\) is the four-root
\(u\)--\(v\) path \(k\)-th, and reading the whole engine zoo off the
constraints.

\textbf{Static-counting path (\(\lv2\)--\(\lv4\), ``the array never changes'').}
\secref{c19-mst} in full: the sorted list per node, the \(\bigO{\log^2 n}\)
count, the external-binary-search \(k\)-th, the index\(\times\)value view.
Then \secref{c19-chairman} for the faster \(k\)-th (a joint descent, no outer
binary search), and \secref{c19-sparse} for the \(O(1)\) idempotent special
case.

\textbf{Persistence path (\(\lv3\)--\(\lv5\), ``read the past'').}
\secref{c19-copy}: path copying, the node pool, the persistent array then the
persistent segment tree, and the fat-node contrast; the hand trace of \(A[3]=9\)
copies five nodes and leaves \(v_0\) intact.  Then \secref{c19-chairman} -- the
prefix versions are persistence with the versions \emph{chosen for you} -- and
\secref{c19-ontrees}, where the versions walk a tree.

\textbf{On-trees path (\(\lv4\)--\(\lv5\), ``the query is a path'').}
\secref{c19-chairman} for the two-root \(k\)-th mechanism, then
\secref{c19-ontrees}: root per DFS node plus Ch.11's LCA turns a \(u\)--\(v\)
path into four roots (SPOJ~COT).  Path \emph{updates} on a tree are the Ch.22
HLD boundary -- forward-cited, no chain code here.
\end{chaptertheorybox}

\begin{figure}[H]
\centering
\begin{tikzpicture}[diagram, x=9mm, y=7mm]
  \foreach \st/\lab in {1/lv1,2/lv2,3/lv3,4/lv4,5/lv5,6/lv6} {
    \fill[\lab!22, draw=\lab, line width=0.8pt] (\st-1,0) rectangle (\st,\st*0.52);
    \node[font=\scriptsize\bfseries, text=\lab] at (\st-0.5,\st*0.52-0.42) {L\st};}
  \node[annot, anchor=south west, font=\scriptsize]
    at (0.05,0.62) {21.1 sorted list per node; one bisection counts values $\le x$};
  \node[annot, anchor=south west, font=\scriptsize]
    at (1.05,1.14) {21.1 build O(n log n), count O(log n) per query; solved SPOJ KQUERY};
  \node[annot, anchor=south west, font=\scriptsize]
    at (2.05,1.66) {21.2 path copying + node pool; keep every version (CF 707D)};
  \node[annot, anchor=south west, font=\scriptsize]
    at (3.05,2.18) {21.3 chairman prefix-versions; k-th by two-root descend (SPOJ MKTHNUM)};
  \node[annot, anchor=south west, font=\scriptsize, text=lv5]
    at (4.05,2.70) {21.4 root per DFS node; k-th on a u-v path from four roots (SPOJ COT)};
  \node[annot, anchor=south west, font=\scriptsize, text=lv6]
    at (5.05,3.22) {21.5 sparse table O(1) RMQ + the whole engine map; Summit};
\end{tikzpicture}
\caption*{Chapter~21 on the six-step ladder: reading a sorted list and
bisecting it is \(\lv2\)--\(\lv3\); a merge-sort-tree count or one hand-built
path-copy is \(\lv3\); the node-pool persistence and the chairman two-root
\(k\)-th are \(\lv4\); the four-root path \(k\)-th and reading the engine zoo
off the constraints are \(\lv5\); the sparse table plus the Summit are step~6.
Ground-up path \(=\) steps 2--4; full contest path runs to~6.}
\end{figure}

% ---------------- The running examples ----------------
\subsection*{The Running Example: One Array, Read Two New Ways}

\emph{One} array carries the whole chapter -- the same 16-slot \(A\) as
Chapters~18, 19 and~20, whose sum segment tree is byte-for-byte the grid
Ch.19's opener fixed (\figref{c13-tree}):
\[
  A=(\,3,2,4,5,1,1,5,3,\;7,0,2,4,1,6,3,2\,),\qquad \textstyle\sum A = 49 .
\]
Chapters~19 and~20 stored \emph{one aggregate} per interval node.  This chapter
asks a different question of the \emph{same} tree -- \emph{how many values in
this interval are \(\le x\)?} -- and answers it by hanging one \term{sorted
list} on each node.  Reading \(A\) through that lens is the merge-sort tree;
then we make one genuinely new move on it, the write \(A[3]=9\) done
\emph{persistently}, so the untouched \(v_0\) survives.

\begin{figure}[H]
\centering
\begin{tikzpicture}[diagram, scale=0.62, transform shape]
  % ---- root t1 full sorted list ----
  \node[boxnode, primaryblue, minimum size=10mm, font=\scriptsize] (t1) at (8,5.6) {\(t_1\)};
  \node[font=\tiny\ttfamily, text=darkgray, anchor=north] at (8,5.05) {\(0,1,1,1,2,2,2,3,3,3,4,4,5,5,6,7\)};
  % ---- level 2: t2, t3 ----
  \node[boxnode, primaryblue, minimum size=9mm, font=\scriptsize] (t2) at (4,3.8) {\(t_2\)};
  \node[font=\tiny\ttfamily, text=darkgray, anchor=north] at (4,3.25) {\(1,1,2,3,3,4,5,5\)};
  \node[boxnode, primaryblue, minimum size=9mm, font=\scriptsize] (t3) at (12,3.8) {\(t_3\)};
  \node[font=\tiny\ttfamily, text=darkgray, anchor=north] at (12,3.25) {\(0,1,2,2,3,4,6,7\)};
  % ---- level 3: t4..t7 ----
  \node[boxnode, successgreen, minimum size=8mm, font=\scriptsize] (t4) at (2,2) {\(t_4\)};
  \node[font=\tiny\ttfamily, text=green!45!black, anchor=north] at (2,1.5) {\(2,3,4,5\)};
  \node[boxnode, successgreen, minimum size=8mm, font=\scriptsize] (t5) at (6,2) {\(t_5\)};
  \node[font=\tiny\ttfamily, text=green!45!black, anchor=north] at (6,1.5) {\(1,1,3,5\)};
  \node[boxnode, primaryblue, minimum size=8mm, font=\scriptsize] (t6) at (10,2) {\(t_6\)};
  \node[font=\tiny\ttfamily, text=darkgray, anchor=north] at (10,1.5) {\(0,2,4,7\)};
  \node[boxnode, primaryblue, minimum size=8mm, font=\scriptsize] (t7) at (14,2) {\(t_7\)};
  \node[font=\tiny\ttfamily, text=darkgray, anchor=north] at (14,1.5) {\(1,2,3,6\)};
  % ---- interval captions ----
  \node[font=\tiny, text=darkgray, anchor=south] at (2,2.55) {\([1..4]\)};
  \node[font=\tiny, text=darkgray, anchor=south] at (6,2.55) {\([5..8]\)};
  \node[font=\tiny, text=darkgray, anchor=south] at (10,2.55) {\([9..12]\)};
  \node[font=\tiny, text=darkgray, anchor=south] at (14,2.55) {\([13..16]\)};
  \node[font=\tiny, text=darkgray, anchor=south] at (4,4.35) {\([1..8]\)};
  \node[font=\tiny, text=darkgray, anchor=south] at (12,4.35) {\([9..16]\)};
  % ---- edges ----
  \draw[flow] (t1) -- (t2); \draw[flow] (t1) -- (t3);
  \draw[flow] (t2) -- (t4); \draw[flow] (t2) -- (t5);
  \draw[flow] (t3) -- (t6); \draw[flow] (t3) -- (t7);
  \node[font=\scriptsize, text=darkgray, anchor=east] at (-0.2,5.6) {root};
\end{tikzpicture}
\caption{The merge-sort tree over \(A\): the \emph{same} heap intervals as
Ch.19's sum tree (\figref{c13-tree}), but each node now stores its interval's
\term{sorted list} instead of a single aggregate.  Only the top three levels are
drawn; a deeper node's list is a contiguous run of the list above it, and the
leaves \(t_{16..31}\) are the single elements.  The two green nodes \(t_4,t_5\)
are the ones the running count query touches.  Build merges lists on the way
up (\(\bigO{n\log n}\)); counting values \(\le x\) over a range bisections each
of the \(O(\log n)\) canonical nodes (\(\bigO{\log^2 n}\)).}
\label{fig:c19-mst}
\end{figure}

\begin{tracebox}[The merge-sort-tree count in full: how many of A(2..9) are at most 3?]
\textbf{Query} \([2..9]\) decomposes into the canonical nodes visited by
\(\coderef{c19-mst-count}\)'s recursion.  Descending from \(t_1[1..16]\): the
range straddles \(t_2[1..8]\) (partially) and \(t_3[9..16]\) (just index~9).
The covered pieces are the leaf \([2..2]=\{2\}\), the node \(t_9[3..4]=\{4,5\}\),
the node \(t_5[5..8]=\{1,1,3,5\}\), and the leaf \([9..9]=\{7\}\).  Each sorted
list is bisected for \(\le 3\):
\[
  \{2\}\!\to 1,\qquad \{4,5\}\!\to 0,\qquad \{1,1,3,5\}\!\to 3,\qquad \{7\}\!\to 0,
\]
so the count is \(1+0+3+0=\boxed{4}\).  Cross-check by hand: \(A[2..9]=
(2,4,5,1,1,5,3,7)\) has exactly \(2,1,1,3\) at most~3 -- four values.  Four
bisections, no re-sorting: the sorted lists were paid for once at build time.
\end{tracebox}

Now the \emph{persistence} act.  We keep \(v_0\) as Ch.19's sum tree, and write
\(A[3]=9\) (old value \(4\), so \(\Delta=+5\)) \emph{without overwriting}: copy
only the five nodes on the root-to-leaf path and share the other twenty-six.

\begin{center}
\scriptsize
\renewcommand{\arraystretch}{1.1}
\setlength{\tabcolsep}{3pt}
\begin{tabular}{@{}l|ccccccccccccccc@{}}
\toprule
\textbf{node} & \(t_1\)&\(t_2\)&\(t_3\)&\(t_4\) & \(t_5\)&\(t_6\)&\(t_7\)&\(t_8\) & \(t_9\)&\(t_{10}\)&\(t_{11}\)&\(t_{12}\)&\(t_{13}\)&\(t_{14}\)&\(t_{15}\)\\
\midrule
\textbf{interval} & 1--16&1--8&9--16&1--4 & 5--8&9--12&13--16&1--2 & 3--4&5--6&7--8&9--10&11--12&13--14&15--16\\
\(v_0\) (\(\sum\)) & 49&24&25&14 & 10&13&12&5 & 9&2&8&7&6&7&5\\
\(v_1\) (\(A_3{=}9\)) & \textbf{54}&\textbf{29}&25&\textbf{19} & 10&13&12&5 & \textbf{14}&2&8&7&6&7&5\\
\bottomrule
\end{tabular}
\end{center}

\noindent{\footnotesize Bold cells are the five \emph{copied} nodes of version
\(v_1\); the leaf \(t_{18}=9\) (and its new sum) is off this internal table.
Every non-bold node is \emph{shared} with \(v_0\) -- \(v_0\) is still a perfect,
readable tree.}

\begin{insightbox}[Persistence is ``copy the touched path, share the rest'' -- and that is five nodes, not thirty-one]
The write \(A[3]=9\) touched the leaf \(t_{18}\) and its four ancestors
\(t_9,t_4,t_2,t_1\) -- exactly \(\log_2 16 + 1 = 5\) of the 31 nodes.  A
\emph{mutable} tree overwrites those five and loses \(v_0\); a
\term{persistent} tree allocates five fresh nodes, points \(v_1\) at the new
root, and leaves every off-path pointer untouched -- so \(v_0\) and \(v_1\) are
two \emph{different trees} that happen to share twenty-six nodes.  Chapter~11
did exactly this on a trie (\defref{c08-persistent}: depth-many copied nodes);
here it is the segment tree, log-many.  The whole cost is the \emph{pool} of
fresh nodes, which \(\secref{c19-copy}\) sizes precisely -- and the reason lazy
tags from Ch.20 sit so badly beside it (a pending tag lives in no single
version).
\end{insightbox}

The second running object is the tree \(T\) Chapters~12--17 have been sharing
-- nodes \(1..10\), root~1, edges
\[
  1\!-\!2,\; 1\!-\!3,\; 2\!-\!4,\; 2\!-\!5,\; 3\!-\!6,\;
  3\!-\!7,\; 5\!-\!8,\; 7\!-\!9,\; 9\!-\!10 ,
\]
with the Ch.16 preorder (\emph{inclusive}):
\[
  \tin=(1,2,6,3,4,7,8,5,9,10),\qquad
  \tout=(10,5,10,3,5,7,10,5,10,10).
\]
For \(\secref{c19-ontrees}\) we hang a weight on each node,
\(w=(3,2,4,5,1,1,5,3,7,0)\) for nodes \(1..10\).  Then the path \(8\!-\!10\) is
\(8,5,2,1,3,7,9,10\) with weights \((3,1,2,3,4,5,7,0)\), whose sorted order is
\((0,1,2,3,3,4,5,7)\) -- so its \(6\)-th smallest is \(4\); the path \(3\!-\!8\)
is \(3,1,2,5,8\) with weights \((4,3,2,1,3)\) sorted \((1,2,3,3,4)\).  The four
roots at the nodes' persistent versions reproduce exactly these multisets, as
\(\secref{c19-ontrees}\) traces.

Three angles, three hand-computed overlays -- each lives in the section named on
its row and never changes:

\begin{center}
\small
\renewcommand{\arraystretch}{1.25}
\begin{tabular}{@{}p{0.22\linewidth} p{0.52\linewidth} l@{}}
\toprule
\textbf{Angle} & \textbf{What comes out} & \textbf{Traced in} \\
\midrule
sorted lists & count \(\le 3\) on \(A[2..9]\): four canonical bisections sum to \(4\) & \secref{c19-mst} \\
path copy & \(A_3{=}9\) copies \(t_1,t_2,t_4,t_9,t_{18}\) (5 of 31); \(v_0\) survives intact & \secref{c19-copy} \\
four roots & \(k\)-th on tree path \(8\!-\!10\) (\(k{=}6\)) \(=4\) from roots \(8,10,\mathit{lca}{=}1\) & \secref{c19-ontrees} \\
\bottomrule
\end{tabular}
\end{center}

% ---------------- Which engine, cost chart ----------------
\subsection*{Which Engine? (the question-shape chart)}

Every structure here answers range questions on \(A\); the only things that
change are \emph{what a node stores} and \emph{what you may still do to the
array}.  That choice is forced by the question and the update model -- and
decisively by whether the array changes at all.

\begin{center}
\small
\renewcommand{\arraystretch}{1.25}
\begin{tabular}{@{}p{0.22\linewidth} p{0.20\linewidth} p{0.20\linewidth} p{0.26\linewidth}@{}}
\toprule
\textbf{Question} & \textbf{Array changes?} & \textbf{Engine} & \textbf{Why this engine} \\
\midrule
count / \(k\)-th values \(\le x\) in \([l,r]\) & never (static) & merge-sort tree & one bisection per canonical node; \(O(\log^2 n)\) count, no updates needed \\
\(k\)-th smallest in \([l,r]\), fast & never & chairman tree & two prefix roots descend jointly, \(O(\log V)\) -- beats the \(\log^3 n\) pipeline \\
read any PAST version after writes & yes, but keep history & persistent segtree & path copying preserves every version; \(n+q\log n\) nodes \\
\(k\)-th / count on a tree \(u\)--\(v\) path & never (values fixed) & persistent roots on \(T\) & root per DFS node + LCA = four-root inclusion--exclusion \\
static range min/max/gcd, \(O(1)\) & never & sparse table & overlap two idempotent blocks; only for idempotent merges (min/max/gcd), \emph{not} sum \\
range \emph{update} + range query, present only & yes, overwrite & lazy segtree (Ch.20) & \(\secref{c19-copy}\) explains why a lazy tag will not persist \\
\bottomrule
\end{tabular}
\end{center}

\begin{insightbox}[One array; the whole game is ``does it change, and do I read the past'']
Rows~1--2 and~5 are \emph{static} -- the array is fixed, so you may precompute
freely (sorted lists, prefix roots, sparse blocks).  Row~3 is the \emph{dynamic
but historical} case -- writes happen but you must still read old versions, so
you path-copy.  Row~6 is the case this chapter deliberately does \emph{not}
serve: range updates you never look back on belong to Ch.20's lazy tag.  When a
persistent problem feels hard, you are choosing what a version's root points at
and how many nodes a write must copy -- not inventing a new traversal.
\end{insightbox}

% ---------------- Taxonomy chart ----------------
\begin{figure}[H]
\centering
\begin{tikzpicture}[diagram, scale=0.95, transform shape,
    mnode/.style={circle, draw=#1, line width=0.9pt, fill=#1!8, minimum size=7mm,
                  font=\scriptsize\bfseries, text=#1!60!black}]
  \node[mnode=primaryblue] (a) at (0,0)     {sorted};
  \node[annot, font=\scriptsize, anchor=west, text width=9.5cm, align=left] at (1.35,0) {each node stores its interval's sorted list; count $\le x$ in a range is one bisection per canonical node, $O(\log^2 n)$ (21.1)};
  \node[mnode=successgreen] (b) at (0,-1)  {copy};
  \node[annot, font=\scriptsize, anchor=west, text width=9.5cm, align=left] at (1.35,-1) {a write copies only the root-to-leaf path (log-many nodes) and shares the rest; the old root is still a full tree (21.2)};
  \node[mnode=difforange]   (c) at (0,-2)  {pool};
  \node[annot, font=\scriptsize, anchor=west, text width=9.5cm, align=left] at (1.35,-2) {fresh nodes come from an explicit pool; budget n + q log n, because the usual 4n array dies the moment versions stack up (21.2)};
  \node[mnode=warningred]   (d) at (0,-3)  {prefix};
  \node[annot, font=\scriptsize, anchor=west, text width=9.5cm, align=left] at (1.35,-3) {$r_i$ = frequency tree over the value axis after $A[1..i]$; count $\le x$ in $[l,r]$ = $r_r$ minus $r_{l-1}$ on matching nodes (21.3)};
  \node[mnode=softpurple]   (e) at (0,-4)  {k-th};
  \node[annot, font=\scriptsize, anchor=west, text width=9.5cm, align=left] at (1.35,-4) {the k-th smallest in $[l,r]$ is a JOINT descent of two roots: left-child counts decide the branch, $O(\log V)$; on a tree path use four roots + LCA (21.3, 21.4)};
  \node[mnode=deepblue]     (f) at (0,-5)  {overlap};
  \node[annot, font=\scriptsize, anchor=west, text width=9.5cm, align=left] at (1.35,-5) {sparse table: st[j][i] folds $2^j$ slots; answer $[l,r]$ by overlapping two blocks -- legal only for idempotent merges (min/max/gcd), not sum (21.5)};
\end{tikzpicture}
\caption{The chapter's six-state taxonomy, read top to bottom as an escalation.
Rows~1--2 are the two static ideas (a merge-sort tree counts, path copying
preserves); row~3 is the memory law that makes persistence actually run; row~4
marries persistence to the value axis (prefix versions); row~5 is the payoff --
\(k\)-th by joint descent, extended to tree paths by four roots; row~6 is the
\(O(1)\) static idempotent special case, the sparse table.  Every solved problem
sits on exactly one row: SPOJ~KQUERY on row~1, CF~707D on rows~2--3, SPOJ
MKTHNUM on rows~4--5, SPOJ~COT on row~5, and the Summit spanning rows~1 and~4
(a \(k\)-th-in-a-range the merge-sort tree samples and the chairman computes
deterministically).}
\label{fig:c19-taxonomy}
\end{figure}

% ---------------- Who owns what ----------------
\subsection*{Who Owns What (so nothing is learned twice)}

\begin{center}
\small
\renewcommand{\arraystretch}{1.25}
\begin{tabular}{@{}p{0.30\linewidth} p{0.30\linewidth} p{0.34\linewidth}@{}}
\toprule
\textbf{Idea} & \textbf{Owner} & \textbf{This chapter's share} \\
\midrule
build / point update / range query, arbitrary combine, iterative \(2n\) tree & Ch.~19 (\chref{c13}) & the interval geometry is ASSUMED; merge-sort and persistent trees reuse the same heap layout -- no re-teach of the plain tree \\
lazy pending tag, range update / range query & Ch.~20 (\chref{c15}) & 21.2 explains only WHY a tag will not persist (it lives in no single version), back-citing \secref{c15-strains}; no lazy code re-appears \\
persistent binary TRIE (\defref{c08-persistent}, LC~1707) & Ch.~11 (\chref{c08}) & the depth-many special case, cited by label; this chapter GENERALISES to whole structures and never re-builds the trie \\
binary lifting / \(\mathit{lca}()\) & Ch.~11 (\chref{c11}) & 21.4 calls \(\mathit{lca}()\) for the four-root path; not re-derived \\
Euler tour, \(\tin/\tout\), ``subtree \(=\) interval'' & Ch.~5 (\chref{c24}) & cited by label in 21.4; this chapter's tree roots are per-DFS-node \\
BIT + offline sweeps (last/second-last position) & Ch.~18 (\chref{c14}) & 21.3 contrasts prefix-versions vs the offline right-end sweep; CF~1000F/813E re-treated in prose, never re-linked \\
heavy-light path queries & Ch.~22 (\chref{c16}) & path \emph{updates} belong there; 21.4 handles \emph{static} path \(k\)-th only \\
segment-tree beats, Li Chao & Ch.~23 (\chref{c29}) & forward-cited as the non-composable / non-idempotent boundary; no beats code here \\
Mo's algorithm (offline) & Ch.~27 (\chref{c21}) & forward-cited as the offline rival to persistence for some range queries; no Mo code here \\
rollback DSU / partial persistence by rollback & Ch.~17 (\chref{c28}) & 21.2 prose-contrasts rollback (one version, undo) vs path copying (many versions, keep); no DSU here \\
\bottomrule
\end{tabular}
\end{center}

% ---------------- Recognition guide ----------------
\begin{recognitionbox}[Pattern card -- ``count or k-th in a range, or read a past version'']
\reckey{``how many of \(a[l..r]\) are \(\le x\) / in \([p,q]\)'', ``the \(k\)-th
smallest in \(a[l..r]\)'' with \emph{no updates} to the array -- a merge-sort
tree or a chairman tree; ``answer queries on the array as it was after the
\(t\)-th update'', ``keep every version'', ``persistent'' -- path copying;
``the \(k\)-th smallest weight on the path \(u\) to \(v\)'' in a fixed tree --
persistent roots plus LCA; and ``range minimum / maximum, many queries, array
never changes, need \(O(1)\)'' -- a sparse table.}
\recuse{
  \item STATIC array + ``count / \(k\)-th values \(\le x\) in a range'' --
        merge-sort tree, \secref{c19-mst} (SPOJ~KQUERY);
  \item ``\(k\)-th smallest in \([l,r]\)'' wanting \(O(\log V)\) -- chairman
        prefix-version tree, \secref{c19-chairman} (SPOJ~MKTHNUM);
  \item ``read the array as it was at version \(t\)'' -- persistent segment tree
        with a node pool, \secref{c19-copy} (CF~707D);
  \item ``\(k\)-th / count on a tree \emph{path}'' -- root per DFS node + LCA,
        \secref{c19-ontrees} (SPOJ~COT);
  \item pure STATIC idempotent RMQ (min/max/gcd), \(O(1)\) -- sparse table,
        \secref{c19-sparse}.
}
\recnotuse{
  \item the array is \emph{updated} and you only ever need the \emph{present}
        sum/min -- that is Ch.~19's plain tree (point writes) or Ch.~20's lazy
        tag (range writes); persistence costs a pool you do not need;
  \item the update is invertible and reads are sums -- Ch.~18's Fenwick tree,
        \chref{c14}; no version is kept;
  \item the query is idempotent-min but the array changes -- a sparse table is
        \emph{static}; go back to Ch.~19/Ch.~20;
  \item queries are offline and unsorted-by-version -- Mo's algorithm,
        Ch.~27 (\chref{c21}), often beats persistence there;
  \item path \emph{updates} on a tree -- heavy-light decomposition, Ch.~22
        (\chref{c16}); 21.4 does static path \(k\)-th only.
}
One-line test: \emph{a range is counted \emph{by value}, or a \emph{past
version} must survive, or a \emph{static} idempotent RMQ wants \(O(1)\) -- then
the sorted / copied / overlapped structure below is what the plain tree cannot
do.}
\end{recognitionbox}

% ---------------- Constraint -> tool ----------------
\subsection*{Reading the Constraints (which engine the limits want)}

\begin{center}
\small
\renewcommand{\arraystretch}{1.25}
\begin{tabular}{@{}p{0.32\linewidth} p{0.40\linewidth} p{0.22\linewidth}@{}}
\toprule
\textbf{Constraint signature} & \textbf{What it rules in/out} & \textbf{Tool} \\
\midrule
static array, \(n\le3\cdot10^4\)/\(q\le2\cdot10^5\), ``count \(\le x\) in \([l,r]\)'' & merge-sort tree \(O(\log^2 n)\)/query; no version, no update & \secref{c19-mst} (SPOJ KQUERY) \\
\(k\)-th smallest in \([l,r]\), \(n\le10^5\), static & chairman tree \(O(\log V)\) beats the merge-sort \(O(\log^3 n)\) pipeline & \secref{c19-chairman} (SPOJ MKTHNUM) \\
\(n,m\le10^3\), \(q\le10^5\) writes \emph{and} ``state after operation \(k\)'' & persistent segtree over a node pool; budget \(n+q\log n\) & \secref{c19-copy} (CF 707D) \\
tree \(N\le10^5\), \(k\)-th weight on a path \(u\)--\(v\) & root per DFS node + LCA = four roots, \(O(\log N)\) & \secref{c19-ontrees} (SPOJ COT) \\
\(n,q\le2\cdot10^5\) static range min/max, no updates & sparse table \(O(1)\)/query; \emph{idempotent} only (not sum) & \secref{c19-sparse} \\
range \emph{update} interleaved with range query, present-only & NOT here -- Ch.~20's lazy tag; persistence cannot share a tag & \secref{c19-copy} boundary \\
\bottomrule
\end{tabular}
\end{center}

% ---------------- Notation ----------------
\subsection*{Notation Used Only Here}

A segment-tree node is named by its heap index \(t_v\) (root \(t_1\)) owning a
closed interval \([l_v,r_v]\), split at \(\mathit{mid}=\lfloor(l_v{+}r_v)/2\rfloor\)
into children \(2v,2v{+}1\); the leaves are \(t_{16..31}=A[1..16]\) exactly as
in Ch.~19.  In a \term{merge-sort tree} \(M_v\) is the \emph{sorted list} stored
at node \(v\), and \(\mathit{lower}(M_v,x)\) is the bisection count of entries
\(\le x\).  In a \term{persistent} tree each version is a root pointer
\(r_0,r_1,\dots\); \code{newnode} bumps a global \(\mathit{pool}\) counter and
returns a fresh index (there is no fixed \([4n]\) array); a node stores
\((\mathit{sum},\mathit{left},\mathit{right})\) where \(\mathit{left},\mathit{right}\)
are \emph{pool indices}, and \term{path copying} clones only the visited path.
A \term{chairman tree} indexes versions by array prefix: \(r_i\) is the
frequency tree over the \emph{value axis} after inserting \(A[1..i]\), and
\([l,r]\) counts are read as \(r_r - r_{l-1}\) on matching nodes; the value
axis is coordinate-compressed to \([1..V]\) when values are large.  On a tree,
\(\mathit{root}_u\) is the persistent root at DFS node \(u\) (holding the
root-to-\(u\) multiset), \(\mathit{par}(\cdot)\) is the parent, and the path
aggregate is \(\mathit{root}_u+\mathit{root}_v-\mathit{root}_w-\mathit{root}_{\mathit{par}(w)}\)
for \(w=\mathit{lca}(u,v)\).  The \term{sparse table} \(\mathit{st}[j][i]\)
folds \(2^j\) elements from position \(i\); \(k=\lfloor\log_2(r-l+1)\rfloor\)
gives the two overlapping blocks.  \(\tin/\tout/\mathit{lca}\) are Ch.~5/Ch.~11's,
used by reference.  Every section runs on \(A\) and \(T\) above: if a table here
and a trace there disagree, this opener is authoritative -- fix the section,
never the other way.
